\subsection{APPLICATION: APPROXIMATION OF CONTINUOUS REAL-VALUED FUNCTIONS DEFINED ON A PRODUCT OF COMPACT SPACES}

THEOREM 4. Let \((\mathbf{X}_t)_{t\in \mathbf{I}}\) be a family of compact spaces, and let

\[
\mathrm{X} = \prod_ {t \in \mathrm{I}} \mathrm{X} _ {t}.
\]

Then every continuous real-valued function on X can be uniformly approximated by sums of a finite number of functions of the form

\[
(x _ {t}) \rightarrow \prod_ {\alpha \in \mathbf {J}} u _ {\alpha} (x _ {\alpha}),
\]

where J is an (arbitrary) finite subset of I and \(u_{\alpha}\) is a continuous real-valued function on \(X_{\alpha}\) for each \(\alpha \in J\) .

Consider the set H of ``functions of one variable'' \((x_{t}) \rightarrow u_{\alpha}(x_{\alpha})\) (any \(\alpha \in I\) ) which are continuous on X. This set separates the points of X, for if \(x = (x_{t})\) and \(y = (y_{t})\) are any two distinct points of X, there exists \(\alpha \in I\) such that \(x_{\alpha} \neq y_{\alpha}\) and there exists a continuous real-valued function \(h_{\alpha}\) on \(X_{\alpha}\) such that \(h_{\alpha}(x_{\alpha}) \neq h_{\alpha}(y_{\alpha})\) . The function \(x \rightarrow h_{\alpha}(\operatorname{pr}_{\alpha} x)\) then belongs to H and takes distinct values at x and y. Since every polynomial in the functions of H is of the form stated in the theorem, the result follows from Theorem 3.

If not all the \(X_{t}\) are compact, the conclusion of Theorem 4 is not necessarily valid (cf.~Exercise 9).

